Написать функцию $f(x)$, вычисляющую и возвращающую куб числа $x$. С ее помощью вычислить кубы чисел $A$, $B$, $C$ и $D$.
Написать функцию $f(x, y)$, вычисляющую и возвращающую сумму среднего арифметического и среднего геометрического чисел $x$ и $y$. Продемонстрировать ее работу на двух парах чисел $A$ и $B$, $C$ и $D$.
Написать функцию $f(a, h)$, вычисляющую и возвращающую площадь треугольника по заданному основанию $a$ и высоте $h$. Продемонстрировать ее работу для двух треугольников.
Написать функцию $f(a, b)$, вычисляющую и возвращающую гипотенузу прямоугольного треугольника по заданным катетам $a$ и $b$. Продемонстрировать ее работу для трех треугольников.
Написать функцию $f(a, b)$, вычисляющую и возвращающую площадь прямоугольника по заданным длинам его сторон $a$ и $b$. Продемонстрировать ее работу для трех прямоугольников.
Написать функцию $f(a, b)$, вычисляющую и возвращающую величину неизвестного угла $C$ в треугольнике по заданным величинам его известных углов $a$ и $b$. Продемонстрировать ее работу для трех треугольников.
Написать функцию $f(a, b)$, вычисляющую и возвращающую длину средней линии трапеции по заданным длинам оснований $a$ и $b$. Продемонстрировать ее работу для трех трапеций.
Написать функцию $f(x, C)$, вычисляющую и возвращающую число $Y$, полученное из целого числа $x$ приписыванием к нему справа цифры $C$. Продемонстрировать ее работу для трех чисел.
Написать функцию $f(x, C)$, вычисляющую и возвращающую число $Y$, полученное из целого числа $x$ приписыванием к нему слева цифры $C$. Продемонстрировать ее работу для трех чисел.
Написать функцию $\text{sign}(x)$. $\text{sign}(x) = 1$, если $x > 0$, $\text{sign}(x) = 0$ во всех остальных случаях. Вычислить для чисел $A$, $B$, $C$: $\text{sign}(A) + \text{sign}(B) + \text{sign}(C)$.
Написать функцию $n(x)$, значение которой равно $0$, если $x > 0$, и единице во всех остальных случаях. Вычислить для чисел $A$, $B$, $C$: $n(A) + n(B) + n(C)$.
Написать функцию $f(a, b, c)$, вычисляющую и возвращающую количество корней квадратного уравнения с коэффициентами $a$, $b$ и $c$.
Написать функцию $f(r)$, вычисляющую и возвращающую площадь круга по значению радиуса $r$. Продемонстрировать работу функции для двух значений радиуса.
Написать функцию $f(r_1, r_2)$, вычисляющую и возвращающую площадь кольца, заключенного между кругами двух радиусов $r_1$ и $r_2$. Продемонстрировать работу функции для двух колец.
Написать функцию $f(r)$, вычисляющую и возвращающую длину окружности по значению радиуса $r$. Продемонстрировать работу функции для двух значений радиуса.
Написать функцию $f(a, h)$, вычисляющую и возвращающую периметр равнобедренного треугольника по его основанию $a$ и высоте $h$. Продемонстрировать ее работу для трех треугольников.
Написать функцию $f(r)$, вычисляющую и возвращающую объем шара по радиусу $r$. Продемонстрировать работу функции для двух значений радиуса.
Написать функцию $f(A, B)$, вычисляющую и возвращающую сумму целых чисел от $A$ до $B$.
Написать функцию $f(A, B, C)$, вычисляющую и возвращающую для чисел $A$ и $B$ их сумму, если $C = 1$, и разность, если $C$ отлично от $1$. Продемонстрировать работу функции для пяти вариантов параметров.
Написать функцию $f(x, y)$, вычисляющую и возвращающую номер четверти на координатной плоскости по координатам точки $(x, y)$. Продемонстрировать работу функции для четырех вариантов параметров, соответствующих всем четвертям.
Написать функцию $f(x, y)$, вычисляющую, входит ли начало координат в отрезок с заданными координатами и лежащий на оси $OX$. Если входит, то функция должна вернуть $1$, иначе --- $0$.
Написать функцию $f(x)$, проверяющую, четное ли число $x$. Если да, то вернуть $1$, иначе --- $0$. Проверить работу функции на $3$ числах.
Написать функцию $f(x)$, проверяющую, является ли число $x$ квадратом другого числа.
Написать функцию $f(x)$, проверяющую, является ли число $x$ простым. Если да, то вернуть $1$, иначе --- $0$. Число называется простым, если имеет два делителя --- $1$ и само себя.
Написать функцию $f(x)$, проверяющую, является ли число $x$ степенью числа $5$. Если да, то вернуть $1$, иначе --- $0$.
Написать функцию $f(a, b, c)$, которая определяет и возвращает максимальное из трех целых чисел, полученных в качестве аргументов.
Написать функцию $f(a, b, c)$, которая определяет и возвращает минимальное из трех целых чисел, полученных в качестве аргументов.
Написать функцию $f(n)$, которая определяет и возвращает произведение четных цифр натурального числа, полученного в качестве аргумента.
Написать функцию $f(n)$, которая определяет и возвращает сумму нечетных цифр натурального числа, полученного в качестве аргумента.
Написать функцию $f(n)$, которая определяет и возвращает количество четных цифр натурального числа, полученного в качестве аргумента.
